% =====================================================================
\section{Implementation}
\label{sec:impl}
% =====================================================================

The implementation is a Lean~4 package (\texttt{CA}) available at
\url{https://github.com/marcellop71/CA} (repository to be made public
upon acceptance).  Downstream projects depend on it with a single line:

\begin{lstlisting}[style=pseudo]
require ca from git "https://github.com/marcellop71/CA" @ "main"
\end{lstlisting}

\paragraph{End-to-end workflow.}  A user who wants to content-address
their project performs four steps: (1)~add \texttt{require ca} to their
\texttt{lakefile.lean} and run \texttt{lake update}; (2)~annotate
declarations with \texttt{@[publish]} or \texttt{@[open\_point]}
(directly or retroactively via \texttt{attribute [publish] SomeName});
(3)~add \texttt{\#ca\_registry "registry/"} to a Lean file that
imports the annotated modules; (4)~run \texttt{lake build}---the
registry folder is generated automatically during compilation, no
separate tool invocation required.  The generated \texttt{registry/}
folder is committed to Git and pushed; the project is now a node in the
decentralized registry.

\paragraph{CLI.}  For batch processing (e.g., indexing all of Mathlib),
a CLI executable provides \texttt{ca address} (compute hashes, export
JSON manifest and TSV edge list) and \texttt{ca fetch} (store to
Redis).  Redis serves as an interactive lookup layer: given a content
address, it returns all known names, types, and modules in
constant time---a capability the static JSON manifest does not provide
efficiently at scale.

\begin{table}[h]
\centering
\begin{tabular}{@{}ll@{}}
  \toprule
  \textbf{Module} & \textbf{Role} \\
  \midrule
  \texttt{CA.Canonical}           & L0 (pure) and L1 (MetaM) canonicalization \\
  \texttt{CA.SHA256}              & SHA-256 via OpenSSL EVP (C FFI) \\
  \texttt{CA.ExprHash}            & Expr serialization, batch hashing, Merkle DAG \\
  \texttt{CA.Export}              & JSON manifest, TSV edge list, statistics \\
  \texttt{CA.Util}                & Dependency collection, declaration classification \\
  \texttt{CA.Registry.Basic}      & Entry types, persistent extensions, query API \\
  \texttt{CA.Registry.Attributes} & \texttt{@[publish]}, \texttt{@[open\_point]} registration \\
  \texttt{CA.Registry.Generate}   & \texttt{\#ca\_registry} command, JSON output \\
  \bottomrule
\end{tabular}
\caption{Module overview.}
\label{tab:modules}
\end{table}

\paragraph{Build and dependencies.}  The default \texttt{lake build}
compiles the \texttt{CA} library, which requires only OpenSSL.  The CLI
executable (\texttt{lake build ca}) additionally requires hiredis for
Redis-backed lookups.  A Nix flake provides a development shell with the
Lean toolchain and all native dependencies for both Linux and macOS.
The implementation currently targets Lean~4.29.0-rc1 and tracks the
Lean release cycle; it depends on the Batteries
library~\cite{batteries} for \texttt{NameMapExtension}.

\simone{Add end-to-end user walkthrough. Current implementation section is half a page.}


% =====================================================================
